\chapter{Couche 2 — L'identité}

Quatre mécanismes pour qu'un compte ne puisse pas être pris par quelqu'un
d'autre.

% -----------------------------------------------------------------------------
\mesure{Le hachage des mots de passe}{mesure 1 · \fichier{app.py}}

\menace
Une base de données finit par fuiter : sauvegarde oubliée, injection SQL,
disque revendu. Si elle contient les mots de passe en clair, tous les comptes
tombent --- et, comme les gens réutilisent leurs mots de passe, des comptes
ailleurs aussi.

\parade
On ne stocke jamais le mot de passe, mais son \textbf{empreinte}, calculée par
une fonction à sens unique et \emph{lente}. La lenteur est la propriété
recherchée : elle ne gêne pas l'utilisateur qui se connecte une fois, mais elle
ruine l'attaquant qui doit en tester des milliards.

\danslecode
\cmd{generate\_password\_hash} de Werkzeug, qui applique \textbf{PBKDF2-SHA256}
avec un sel aléatoire par compte et un grand nombre d'itérations. Le sel rend
les tables précalculées inutilisables : deux utilisateurs ayant le même mot de
passe ont deux empreintes différentes.

\begin{lstlisting}[language=Python]
user_id = db.create_user(username, generate_password_hash(password), email or None)
\end{lstlisting}

\begin{piege}{Comparer en temps constant, et même quand le compte n'existe pas}
Une vérification naïve répond \emph{plus vite} quand l'identifiant est inconnu
--- il n'y a pas d'empreinte à calculer. En chronométrant les réponses, on
énumère donc les comptes existants sans jamais réussir une connexion.

La parade tient en une ligne : quand l'utilisateur n'existe pas, on calcule
quand même une empreinte factice, pour que les deux cas coûtent le même temps.

\begin{lstlisting}[language=Python]
ok = user is not None and check_password_hash(user["password_hash"], password)
# meme cout quand le compte n'existe pas : on ne doit pas pouvoir les enumerer
check_password_hash(generate_password_hash("x"), password) if user is None else None
\end{lstlisting}
\end{piege}

\limite
Un mot de passe faible reste faible, même bien haché. D'où la mesure suivante.

% -----------------------------------------------------------------------------
\mesure{La politique de mot de passe}{mesure 2 · \fichier{config.py}}

\menace
Les attaques par dictionnaire ne cassent pas la cryptographie : elles essaient
\cmd{123456}, \cmd{password}, \cmd{azertyuiop}. Un seul compte faible suffit à
entrer.

\parade
Une longueur minimale et le refus d'une liste de mots de passe courants, à
l'inscription.

\danslecode
\begin{lstlisting}[language=Python]
PASSWORD_MIN_LENGTH = int(os.environ.get("INSACLOUD_PASSWORD_MIN", "10"))
COMMON_PASSWORDS = {"password", "motdepasse", "123456789", "1234567890",
                    "azertyuiop", ...}
\end{lstlisting}

Dix caractères, pas huit : la règle classique des huit caractères date d'une
époque où le matériel était mille fois plus lent.

\limite
Aucune politique n'empêche \cmd{Soutenance2026!} --- conforme, et devinable par
qui connaît le contexte. La vraie défense contre le devinage ciblé est la
suivante.

% -----------------------------------------------------------------------------
\mesure{La limitation des tentatives}{mesure 3 · \fichier{app.py}}

\menace
Avec un formulaire de connexion public, un robot teste des milliers de
combinaisons par minute. Sans limite, il finit par trouver.

\parade
Un compteur de tentatives échouées, sur \textbf{deux axes} et dans une fenêtre
glissante.

\danslecode
\begin{lstlisting}[language=Python]
LOGIN_WINDOW_MINUTES    = 15   # la fenetre glissante
LOGIN_MAX_FAILURES_USER = 5    # par compte vise
LOGIN_MAX_FAILURES_IP   = 20   # par adresse source
\end{lstlisting}

Les deux axes traitent deux attaques différentes :

\begin{itemize}[leftmargin=14pt,itemsep=2pt]
\item \textbf{par compte} --- l'attaquant s'acharne sur \cmd{admin} depuis mille
      adresses : le compteur du compte le bloque ;
\item \textbf{par adresse} --- l'attaquant essaie un mot de passe très probable
      sur mille comptes différents (\emph{password spraying}), ce qui ne
      déclenche aucun compteur par compte : le compteur par adresse le bloque.
\end{itemize}

Chaque tentative est enregistrée en base avec son horodatage et son adresse
source ; les entrées de plus de vingt-quatre heures sont purgées à chaque
connexion réussie, pour que la table ne grossisse pas indéfiniment.

\begin{retenir}
Deux compteurs valent mieux qu'un : un seul compteur par compte ne voit pas le
\emph{password spraying}, et un seul compteur par adresse ne voit pas l'attaque
distribuée. C'est la combinaison qui couvre les deux.
\end{retenir}

\limite
Un attaquant disposant de beaucoup d'adresses et visant beaucoup de comptes
passe sous les deux seuils. À l'échelle de ce projet c'est acceptable ; en
production on ajouterait un CAPTCHA après quelques échecs, ou une temporisation
croissante.

% -----------------------------------------------------------------------------
\mesure{La validation des entrées}{mesure 4 · \fichier{config.py}}

\menace
Toute donnée venue de l'utilisateur finit quelque part : dans une requête SQL,
dans un nom de conteneur, dans une page HTML, dans une ligne de commande. Si
elle n'est pas contrainte, elle y fait ce qu'elle veut.

\parade
Des expressions régulières \textbf{ancrées} --- \cmd{\^{}} et \cmd{\$} --- qui
décrivent ce qui est \emph{autorisé}, jamais ce qui est interdit. Une liste
d'interdits est toujours incomplète ; une liste d'autorisés ne l'est jamais.

\danslecode
\begin{lstlisting}[language=Python]
USERNAME_RE   = re.compile(r"^[a-zA-Z0-9][a-zA-Z0-9_.-]{2,31}$")  # compatible Docker
EMAIL_RE      = re.compile(r"^[^@\s]{1,64}@[^@\s.]+(\.[^@\s.]+)+$")
SSH_PUBKEY_RE = re.compile(r"^(ssh-ed25519|ssh-rsa|ecdsa-sha2-nistp(256|384|521)"
                           r"|sk-ssh-ed25519@openssh\.com) [A-Za-z0-9+/]+=*"
                           r"( [^\r\n]{0,128})?$")
MIN_DURATION_MINUTES = 1
MAX_DURATION_MINUTES = 120
MAX_INSTANCES_PER_USER = 3
\end{lstlisting}

Le format de l'identifiant n'est pas arbitraire : il est choisi pour être
\textbf{compatible avec les noms de conteneurs Docker}, puisqu'il s'y retrouve.
Un identifiant contenant une espace ou un point-virgule serait refusé bien
avant d'atteindre la ligne de commande.

\begin{notion}{La requête est bornée avant même d'être lue}
\cmd{MAX\_CONTENT\_LENGTH = 16 * 1024} : Flask rejette toute requête de plus de
seize kilooctets. Aucun formulaire légitime de la plateforme n'en approche ---
le plus gros est une clé publique SSH, quelques centaines d'octets. Cette
borne coupe court aux tentatives d'épuisement mémoire par un corps de requête
géant, avant que le code applicatif ne voie quoi que ce soit.
\end{notion}

\limite
La validation protège le format, pas le sens. Une durée de 120 minutes est
valide ; demander trois machines de 120 minutes toutes les deux minutes reste
possible, et c'est le quota (\cmd{MAX\_INSTANCES\_PER\_USER}) qui l'arrête.
